\documentclass[14pt,a4paper]{extarticle}
\usepackage[T2A]{fontenc}
\usepackage[utf8]{inputenc}
\usepackage[english,russian]{babel}
\usepackage{lmodern}
\usepackage{geometry}
\usepackage{array}
\usepackage{tabularx}
\usepackage{longtable}
\usepackage{ltablex}
\usepackage{booktabs}
\usepackage{hyperref}
\usepackage{indentfirst}
\usepackage{microtype}
\usepackage{graphicx}
\usepackage{setspace}
\usepackage{listings}
\usepackage{graphicx}
\usepackage{caption}
\usepackage{float}
\usepackage{tocloft}
\usepackage{titlesec}
\usepackage{newtxtext,newtxmath}

\titleformat{\section}        % форматируем \section
  {\normalfont\bfseries}      % стиль: обычный + жирный
  {\thesection.}              % номер секции с точкой
  {2mm}                       % расстояние между номером и текстом
  {}                           % перед заголовком можно добавить что-то
\titlespacing*{\section}{1.5cm}{*3}{*1}  % {отступ слева}{отступ сверху}{отступ снизу}

\titleformat{\subsection}        % форматируем \section
  {\normalfont\bfseries}      % стиль: обычный + жирный
  {\thesubsection.}              % номер секции с точкой
  {2mm}                       % расстояние между номером и текстом
  {}                           % перед заголовком можно добавить что-то
\titlespacing*{\subsection}{1.5cm}{*3}{*1}  % {отступ слева}{отступ сверху}{отступ снизу}

\captionsetup{justification=centering}
\renewcommand{\ttdefault}{cmtt}
\keepXColumns
\setlength{\parindent}{3em}
\newgeometry{left=2.5cm,right=2.5cm,top=3cm,bottom=2.5cm}

\renewcommand{\contentsname}{\hspace*{4em}{\smallСодержание}}
\setlength{\cftbeforesecskip}{0.5em}
\renewcommand{\cftsecleader}{\cftdotfill{\cftdotsep}}
\renewcommand{\cftsecaftersnum}{.\hspace{2mm}}

\title{Отчёт по "Управление ИТ-проектами"}
\author{Выполнил студент группы М-РИТ-25-1, Станиславчук Сергей}

\begin{document}

\setstretch{1.5}
\noindent
\begin{minipage}[t]{0.01\textwidth}
    \includegraphics{~/programmer/++Programmer/5_1/LSTU.jpg}
\end{minipage}%
\begin{minipage}[c]{1.1\linewidth}
    \centering
    {\small
    МИНИСТЕРСТВО НАУКИ И ВЫСШЕГО ОБРАЗОВАНИЯ РФ \\[1em]
    ФЕДЕРАЛЬНОЕ ГОСУДАРСТВЕННОЕ БЮДЖЕТНОЕ \\
    ОБРАЗОВАТЕЛЬНОЕ УЧРЕЖДЕНИЕ ВЫСШЕГО ОБРАЗОВАНИЯ \\
    «ЛИПЕЦКИЙ ГОСУДАРСТВЕННЫЙ ТЕХНИЧЕСКИЙ УНИВЕРСИТЕТ»
    }
\end{minipage}%

\vspace{3em}

\noindent
\hspace{1em}%
\begin{minipage}[t]{0.8\textwidth}
\textbf{Институт} \hspace{2em} компьютерных наук \\
\textbf{Кафедра} \hspace{2em} автоматизированных систем управления
\end{minipage}


\vspace{6em}

\begin{center}
\textbf{Практическая работа №6} \\
по дисциплине \\ <<Управление ИТ-проектами>>
\end{center}

\vspace{6em}

\noindent
\begin{tabular*}{\textwidth}{@{\extracolsep{\fill}} l c r }
    Студент М-РИТ-25-1 & \makebox[4cm]{\hrulefill} & Станиславчук С.М. \\[2em]
Руководитель & \\
К.т.н., доцент & \makebox[4cm]{\hrulefill} & Ведищев В.В. \\
\end{tabular*}

\vfill
\begin{center}
    Липецк, 2026
\end{center}

\thispagestyle{empty}

\newpage
\tableofcontents
\newpage

\section{Основное содержание SWOT-анализа}

SWOT-анализ направлен на выявление внутренних характеристик проекта (сильные и слабые стороны), а также внешних факторов (возможности и угрозы), влияющих на его реализацию. В контексте разрабатываемой системы корпоративного обмена сообщениями данный анализ позволяет оценить её применимость, устойчивость к изменениям внешней среды и уровень конкурентоспособности.

\section{Сильные стороны}

Ключевым преимуществом системы является её децентрализованный характер. Развёртывание внутри корпоративной инфраструктуры исключает зависимость от внешних сервисов и снижает риски блокировок, а также устраняет типичную для централизованных платформ единую точку отказа. Это особенно актуально в условиях ограниченного доступа к внешним коммуникационным сервисам.

Дополнительным фактором является повышенный уровень безопасности. Поскольку все данные остаются внутри организации, существенно упрощается контроль доступа, аудит действий пользователей и соблюдение внутренних политик информационной безопасности. Это позволяет минимизировать риски утечек и несанкционированного доступа.

Существенным преимуществом также является высокая производительность системы. Использование локальной сети обеспечивает минимальные задержки и высокую скорость передачи данных, что критично для голосовой связи и передачи файлов. Кроме того, полный контроль над инфраструктурой позволяет организации самостоятельно управлять обновлениями, конфигурацией и политиками хранения данных.

\section{Слабые стороны, проблемы}

Несмотря на перечисленные преимущества, проект обладает рядом ограничений. В первую очередь, это сложность внедрения, связанная с необходимостью развёртывания серверной части, настройки сетевого взаимодействия и обеспечения отказоустойчивости. Это увеличивает сроки ввода системы в эксплуатацию.

Также существенным фактором является потребность в квалифицированном персонале. Для сопровождения системы требуются специалисты, способные обеспечивать её стабильную работу, обновление и устранение возможных инцидентов.

Ограничением является и ориентация системы на локальную сеть. Для обеспечения доступа из внешней среды необходимо внедрение дополнительных решений, таких как VPN или проксирование, что усложняет архитектуру. Кроме того, отсутствие готовой экосистемы и необходимость самостоятельной реализации многих компонентов увеличивают трудозатраты на разработку и сопровождение системы.

\section{Возможности}

Внешняя среда создаёт ряд условий, способствующих внедрению подобных решений. Рост требований к защите информации стимулирует организации переходить на системы, обеспечивающие полный контроль над данными. В этом контексте локальные решения становятся более предпочтительными по сравнению с облачными сервисами.

Дополнительным фактором является политика импортозамещения, в рамках которой возникает необходимость замены зарубежных продуктов внутренними разработками. Это открывает возможности для внедрения системы в корпоративной среде.

Система также может быть интегрирована с существующей инфраструктурой организации, включая системы аутентификации и внутренние сервисы, что повышает её ценность. В перспективе возможно расширение функциональности до уровня полноценной коммуникационной платформы с поддержкой дополнительных сервисов.

\section{События внешней среды как возможности и как угрозы}

Отдельные факторы внешней среды могут одновременно выступать как возможностями, так и угрозами. В частности, ограничение доступа к централизованным мессенджерам повышает востребованность автономных решений, однако может сопровождаться нестабильностью требований и необходимостью оперативной адаптации системы.

Изменения законодательства также имеют двойственный характер. С одной стороны, они могут стимулировать внедрение защищённых систем, с другой — требовать доработки архитектуры и усложнять процесс эксплуатации.

Развитие технологий создаёт возможности для повышения эффективности системы, однако требует постоянного обновления и сопровождения. Наличие альтернативных решений может ускорять внедрение за счёт заимствования подходов, но одновременно усиливает конкуренцию.

\section{Анализ конкурентных сил}

На рынке уже присутствуют зрелые корпоративные решения, такие как Slack, Microsoft Teams, Mattermost и Rocket.Chat. Эти платформы предлагают широкий функционал, включающий текстовую, голосовую и видеосвязь, а также развитые механизмы интеграции с другими сервисами. Они характеризуются высокой стабильностью и наличием сформированной экосистемы.

В то же время такие решения часто ориентированы на облачную инфраструктуру или требуют сложного развёртывания. Это ограничивает их применимость в условиях жёстких требований к безопасности и независимости от внешних сервисов. В результате формируется ниша для специализированных систем, оптимизированных под использование внутри корпоративной сети.

\section{Использование результатов SWOT-анализа}

Результаты проведённого анализа используются для обоснования архитектурных решений и определения направлений развития проекта. В частности, они позволяют учитывать требования к безопасности и независимости, минимизировать выявленные слабые стороны за счёт упрощения внедрения и автоматизации процессов, а также эффективно использовать внешние возможности, связанные с изменениями в среде.

Таким образом, SWOT-анализ служит инструментом стратегического планирования и позволяет повысить обоснованность принимаемых проектных решений, обеспечивая баланс между функциональностью, безопасностью и затратами на реализацию системы.

\end{document}
